% Options for packages loaded elsewhere
\PassOptionsToPackage{unicode}{hyperref}
\PassOptionsToPackage{hyphens}{url}
\documentclass[
  ignorenonframetext,
]{beamer}
\newif\ifbibliography
\usepackage{pgfpages}
\setbeamertemplate{caption}[numbered]
\setbeamertemplate{caption label separator}{: }
\setbeamercolor{caption name}{fg=normal text.fg}
\beamertemplatenavigationsymbolsempty
% remove section numbering
\setbeamertemplate{part page}{
  \centering
  \begin{beamercolorbox}[sep=16pt,center]{part title}
    \usebeamerfont{part title}\insertpart\par
  \end{beamercolorbox}
}
\setbeamertemplate{section page}{
  \centering
  \begin{beamercolorbox}[sep=12pt,center]{section title}
    \usebeamerfont{section title}\insertsection\par
  \end{beamercolorbox}
}
\setbeamertemplate{subsection page}{
  \centering
  \begin{beamercolorbox}[sep=8pt,center]{subsection title}
    \usebeamerfont{subsection title}\insertsubsection\par
  \end{beamercolorbox}
}
% Prevent slide breaks in the middle of a paragraph
\widowpenalties 1 10000
\raggedbottom
\AtBeginPart{
  \frame{\partpage}
}
\AtBeginSection{
  \ifbibliography
  \else
    \frame{\sectionpage}
  \fi
}
\AtBeginSubsection{
  \frame{\subsectionpage}
}
\usepackage{iftex}
\ifPDFTeX
  \usepackage[T1]{fontenc}
  \usepackage[utf8]{inputenc}
  \usepackage{textcomp} % provide euro and other symbols
\else % if luatex or xetex
  \usepackage{unicode-math} % this also loads fontspec
  \defaultfontfeatures{Scale=MatchLowercase}
  \defaultfontfeatures[\rmfamily]{Ligatures=TeX,Scale=1}
\fi
\usepackage{lmodern}
\ifPDFTeX\else
  % xetex/luatex font selection
\fi
% Use upquote if available, for straight quotes in verbatim environments
\IfFileExists{upquote.sty}{\usepackage{upquote}}{}
\IfFileExists{microtype.sty}{% use microtype if available
  \usepackage[]{microtype}
  \UseMicrotypeSet[protrusion]{basicmath} % disable protrusion for tt fonts
}{}
\makeatletter
\@ifundefined{KOMAClassName}{% if non-KOMA class
  \IfFileExists{parskip.sty}{%
    \usepackage{parskip}
  }{% else
    \setlength{\parindent}{0pt}
    \setlength{\parskip}{6pt plus 2pt minus 1pt}}
}{% if KOMA class
  \KOMAoptions{parskip=half}}
\makeatother
\usepackage{color}
\usepackage{fancyvrb}
\newcommand{\VerbBar}{|}
\newcommand{\VERB}{\Verb[commandchars=\\\{\}]}
\DefineVerbatimEnvironment{Highlighting}{Verbatim}{commandchars=\\\{\}}
% Add ',fontsize=\small' for more characters per line
\newenvironment{Shaded}{}{}
\newcommand{\AlertTok}[1]{\textcolor[rgb]{1.00,0.00,0.00}{\textbf{#1}}}
\newcommand{\AnnotationTok}[1]{\textcolor[rgb]{0.38,0.63,0.69}{\textbf{\textit{#1}}}}
\newcommand{\AttributeTok}[1]{\textcolor[rgb]{0.49,0.56,0.16}{#1}}
\newcommand{\BaseNTok}[1]{\textcolor[rgb]{0.25,0.63,0.44}{#1}}
\newcommand{\BuiltInTok}[1]{\textcolor[rgb]{0.00,0.50,0.00}{#1}}
\newcommand{\CharTok}[1]{\textcolor[rgb]{0.25,0.44,0.63}{#1}}
\newcommand{\CommentTok}[1]{\textcolor[rgb]{0.38,0.63,0.69}{\textit{#1}}}
\newcommand{\CommentVarTok}[1]{\textcolor[rgb]{0.38,0.63,0.69}{\textbf{\textit{#1}}}}
\newcommand{\ConstantTok}[1]{\textcolor[rgb]{0.53,0.00,0.00}{#1}}
\newcommand{\ControlFlowTok}[1]{\textcolor[rgb]{0.00,0.44,0.13}{\textbf{#1}}}
\newcommand{\DataTypeTok}[1]{\textcolor[rgb]{0.56,0.13,0.00}{#1}}
\newcommand{\DecValTok}[1]{\textcolor[rgb]{0.25,0.63,0.44}{#1}}
\newcommand{\DocumentationTok}[1]{\textcolor[rgb]{0.73,0.13,0.13}{\textit{#1}}}
\newcommand{\ErrorTok}[1]{\textcolor[rgb]{1.00,0.00,0.00}{\textbf{#1}}}
\newcommand{\ExtensionTok}[1]{#1}
\newcommand{\FloatTok}[1]{\textcolor[rgb]{0.25,0.63,0.44}{#1}}
\newcommand{\FunctionTok}[1]{\textcolor[rgb]{0.02,0.16,0.49}{#1}}
\newcommand{\ImportTok}[1]{\textcolor[rgb]{0.00,0.50,0.00}{\textbf{#1}}}
\newcommand{\InformationTok}[1]{\textcolor[rgb]{0.38,0.63,0.69}{\textbf{\textit{#1}}}}
\newcommand{\KeywordTok}[1]{\textcolor[rgb]{0.00,0.44,0.13}{\textbf{#1}}}
\newcommand{\NormalTok}[1]{#1}
\newcommand{\OperatorTok}[1]{\textcolor[rgb]{0.40,0.40,0.40}{#1}}
\newcommand{\OtherTok}[1]{\textcolor[rgb]{0.00,0.44,0.13}{#1}}
\newcommand{\PreprocessorTok}[1]{\textcolor[rgb]{0.74,0.48,0.00}{#1}}
\newcommand{\RegionMarkerTok}[1]{#1}
\newcommand{\SpecialCharTok}[1]{\textcolor[rgb]{0.25,0.44,0.63}{#1}}
\newcommand{\SpecialStringTok}[1]{\textcolor[rgb]{0.73,0.40,0.53}{#1}}
\newcommand{\StringTok}[1]{\textcolor[rgb]{0.25,0.44,0.63}{#1}}
\newcommand{\VariableTok}[1]{\textcolor[rgb]{0.10,0.09,0.49}{#1}}
\newcommand{\VerbatimStringTok}[1]{\textcolor[rgb]{0.25,0.44,0.63}{#1}}
\newcommand{\WarningTok}[1]{\textcolor[rgb]{0.38,0.63,0.69}{\textbf{\textit{#1}}}}
\usepackage{graphicx}
\makeatletter
\newsavebox\pandoc@box
\newcommand*\pandocbounded[1]{% scales image to fit in text height/width
  \sbox\pandoc@box{#1}%
  \Gscale@div\@tempa{\textheight}{\dimexpr\ht\pandoc@box+\dp\pandoc@box\relax}%
  \Gscale@div\@tempb{\linewidth}{\wd\pandoc@box}%
  \ifdim\@tempb\p@<\@tempa\p@\let\@tempa\@tempb\fi% select the smaller of both
  \ifdim\@tempa\p@<\p@\scalebox{\@tempa}{\usebox\pandoc@box}%
  \else\usebox{\pandoc@box}%
  \fi%
}
% Set default figure placement to htbp
\def\fps@figure{htbp}
\makeatother
\setlength{\emergencystretch}{3em} % prevent overfull lines
\providecommand{\tightlist}{%
  \setlength{\itemsep}{0pt}\setlength{\parskip}{0pt}}
\usepackage{bookmark}
\IfFileExists{xurl.sty}{\usepackage{xurl}}{} % add URL line breaks if available
\urlstyle{same}
\hypersetup{
  hidelinks,
  pdfcreator={LaTeX via pandoc}}

\author{\texorpdfstring{}{}}
\date{}

\begin{document}

\section{Topos Theory: Classifying Spaces and Morita
Equivalence}\label{sec:topos-theory}

\begin{frame}{From Categories to Toposes: The Translation Problem}
\protect\phantomsection\label{from-categories-to-toposes-the-translation-problem}
The preceding sections have shown how case systems, categorial grammars,
distributional semantics, and enriched categories each provide a
different ``window'' onto the same underlying linguistic phenomenon. But
how do we know that results proved in one framework transfer to another?
This is the problem of \emph{inter-theoretic translation}---and it is
precisely the problem that Caramello's {[}-@caramello2016bridges{]}
topos-theoretic bridge technique was designed to solve.
\end{frame}

\begin{frame}{Classifying Toposes and Morita Equivalence}
\protect\phantomsection\label{classifying-toposes-and-morita-equivalence}
\begin{block}{What Is a Topos? Generalized Universes of Sets}
\protect\phantomsection\label{what-is-a-topos-generalized-universes-of-sets}
A \textbf{topos} is a category that behaves like a generalized universe
of sets: it has products, exponentials, a subobject classifier (playing
the role of a ``truth-value object''), and enough structure to interpret
first-order logic internally. The category \textbf{Set} of ordinary sets
is a topos, but there are many others---presheaf categories, sheaf
categories over topological spaces, and the categories of models of
geometric theories.
\end{block}

\begin{block}{Classifying Toposes of Geometric Theories}
\protect\phantomsection\label{classifying-toposes-of-geometric-theories}
Every \emph{geometric theory} \(\mathbb{T}\) (a theory axiomatized by
sequents involving only finite conjunctions, arbitrary disjunctions, and
existential quantification) has a \textbf{classifying topos}
\(\mathcal{E}_\mathbb{T}\)---a canonical topos whose models correspond
exactly to the geometric functors from \(\mathcal{E}_\mathbb{T}\) to
other toposes. The classifying topos encodes the theory's ``logical
shape'' independently of any particular model.

Caramello's key insight {[}-@caramello2016bridges;
-@caramello2021five{]} is that different theories can share the same
classifying topos up to equivalence---an equivalence known as
\textbf{Morita equivalence}. When two theories \(\mathbb{T}_1\) and
\(\mathbb{T}_2\) are Morita equivalent
(\(\mathcal{E}_{\mathbb{T}_1} \simeq \mathcal{E}_{\mathbb{T}_2}\)), any
property that can be expressed as an invariant of the classifying topos
transfers automatically between them.
\end{block}
\end{frame}

\begin{frame}{Bridge Theorems for the Four Case Theories}
\protect\phantomsection\label{bridge-theorems-for-the-four-case-theories}
\begin{block}{Formalizing Case Theories as Geometric Theories}
\protect\phantomsection\label{formalizing-case-theories-as-geometric-theories}
We formalize each case-theoretic framework as a geometric theory:

\begin{enumerate}
\item
  \textbf{Typological case theory} \(\mathbb{T}_{\text{typ}}\): Objects
  are case roles, morphisms are grammatical relations, axioms specify
  alignment constraints (e.g., ``S = A'' in accusative alignment).
\item
  \textbf{Type-logical case theory} \(\mathbb{T}_{\text{log}}\): Objects
  are syntactic types, morphisms are Lambek calculus derivations, axioms
  specify well-formedness conditions (pregroup contractions).
\item
  \textbf{Distributional case theory} \(\mathbb{T}_{\text{dist}}\):
  Objects are vector spaces, morphisms are linear maps, axioms specify
  the composition law for distributional meaning (DisCoCat).
\item
  \textbf{Enriched case theory} \(\mathbb{T}_{\text{enr}}\): Objects are
  expressions, morphisms carry \([0,1]\)-valued weights, axioms specify
  the identity and composition inequalities.
\end{enumerate}
\end{block}

\begin{block}{The Bridge Theorem: Morita Equivalence Chain}
\protect\phantomsection\label{the-bridge-theorem-morita-equivalence-chain}
The central claim is that these four theories are related by a chain of
Morita equivalences:

\[\mathcal{E}_{\mathbb{T}_{\text{typ}}} \leftarrow \mathcal{E}_{\mathbb{T}_{\text{bridge}}} \rightarrow \mathcal{E}_{\mathbb{T}_{\text{log}}} \leftarrow \mathcal{E}_{\mathbb{T}_{\text{bridge}'}} \rightarrow \mathcal{E}_{\mathbb{T}_{\text{dist}}} \]
\{\#eq:eq-6-1\}

where the intermediate toposes are constructed by finding geometric
theories that are simultaneously interpretable in both flanking
frameworks. The Morita equivalence ensures that:

\begin{itemize}
\tightlist
\item
  \textbf{Syntactic theorems port to semantics}: A commutativity result
  proved in the type-logical framework automatically yields a
  compositionality result in the distributional framework.
\item
  \textbf{Typological universals constrain distributional models}:
  Alignment types (accusative, ergative) impose structural constraints
  on the vector spaces used in DisCoCat.
\item
  \textbf{Enriched structure enriches all frameworks}: The
  \([0,1]\)-valued weights from the enriched theory can be pulled back
  to give probabilistic interpretations of typological, logical, and
  distributional constructions.
\end{itemize}

\autoref{fig:functor-alignment} visualizes the alignment functor between
the Accusative and Ergative category instantiations, showing how the
same eight case roles are connected by different morphism structures
under different alignment types.

\begin{figure}
\centering
\pandocbounded{\includegraphics[keepaspectratio,alt={The alignment functor F\textbackslash colon\textbackslash mathcal\{C\}\_\{\textbackslash text\{acc\}\} \textbackslash to \textbackslash mathcal\{C\}\_\{\textbackslash text\{erg\}\} mapping the Accusative case category (source, blue panel) to the Ergative case category (target, amber panel). All eight case roles appear in both categories; purple horizontal arrows show the functor's object-level mapping F(\textbackslash text\{role\}). The key structural difference is in the morphism grouping: the Accusative category groups \textbackslash\{S, A\textbackslash\} \textbackslash to \textbackslash text\{NOM\} (kernel \textbackslash\{(S,A)\textbackslash\}, cf. {[}@eq:eq-2-3{]}), while the Ergative groups \textbackslash\{S, P\textbackslash\} \textbackslash to \textbackslash text\{ABS\} (kernel \textbackslash\{(S,P)\textbackslash\}, cf. {[}@eq:eq-2-4{]}). The functor preserves the role inventory but restructures the morphism pattern---formalizing the alignment-as-functor principle. This diagram provides one link in the Morita equivalence chain of {[}@eq:eq-6-1{]} that enables inter-theoretic bridge transfer.}]{output/figures/functor_alignment.png}}
\caption{The alignment functor
\(F\colon\mathcal{C}_{\text{acc}} \to \mathcal{C}_{\text{erg}}\) mapping
the Accusative case category (source, blue panel) to the Ergative case
category (target, amber panel). All eight case roles appear in both
categories; purple horizontal arrows show the functor's object-level
mapping \(F(\text{role})\). The key structural difference is in the
\textbf{morphism grouping}: the Accusative category groups
\(\{S, A\} \to \text{NOM}\) (kernel \(\{(S,A)\}\), cf.
{[}@eq:eq-2-3{]}), while the Ergative groups \(\{S, P\} \to \text{ABS}\)
(kernel \(\{(S,P)\}\), cf. {[}@eq:eq-2-4{]}). The functor preserves the
role inventory but restructures the morphism pattern---formalizing the
alignment-as-functor principle. This diagram provides one link in the
Morita equivalence chain of {[}@eq:eq-6-1{]} that enables
inter-theoretic bridge transfer.}\label{fig:functor-alignment}
\end{figure}
\end{block}
\end{frame}

\begin{frame}{Phillips and the Universal Language of Thought}
\protect\phantomsection\label{phillips-and-the-universal-language-of-thought}
Phillips {[}-@phillips2024lot{]} provides a striking application of
topos-theoretic methods to cognitive science. He shows that the
\textbf{Language of Thought} (LoT) hypothesis---the claim that cognition
operates over structured, combinatorial representations with
language-like properties---can be formalized categorically, and that the
resulting structure is \emph{universal} in the topos-theoretic sense.

Specifically, Phillips demonstrates that:

\begin{enumerate}
\item
  LoT properties (discrete constituents, role-filler independence,
  systematicity) arise as \textbf{universal constructions} in a
  topos---categorical products, fiber bundles, and presheaves.
\item
  Every topos supports an internal first-order logic, explaining how
  LoT-like logical capacities can emerge in systems (biological or
  artificial) whose internal architecture forms a topos.
\item
  The ``shape'' of cognitive representations is fundamentally
  \textbf{topological}, captured by presheaves and fiber bundles rather
  than by point-set structures.
\end{enumerate}

For our framework, Phillips's result is significant because it provides
a topos-theoretic foundation for the claim that case structure is a
universal feature of cognitive architecture. If the Language of Thought
is topos-universal, and case categories are definable within any topos
(which they are, being small categories with first-order axioms), then
every cognitive system with LoT-like structure must be able to represent
case distinctions---a strong universality claim that goes beyond mere
typological observation.
\end{frame}

\begin{frame}{Syntactic Learning via Classifying Toposes}
\protect\phantomsection\label{syntactic-learning-via-classifying-toposes}
Caramello {[}-@caramello2023syntactic{]} extends the bridge technique to
a learning theory: she shows that classifying toposes can be used to
\emph{learn} the theory of a mathematical structure from finite data (a
finite set of models). The learning algorithm constructs a classifying
topos from the observed data and then extracts the axioms of the
underlying theory.

Applied to case systems, this suggests a principled approach to
\emph{grammatical induction}: given a corpus annotated with case labels,
one could construct the classifying topos of the implicit case theory
and read off its axioms---recovering the alignment type, the morphism
structure, and the enriched weights from data alone. The procedure
operates in four phases:

\begin{enumerate}
\tightlist
\item
  \textbf{Extraction}: Parse a Universal Dependencies treebank for a
  target language, collecting all case-labeled dependency arcs. Each arc
  \((r_1, \text{rel}, r_2)\) instantiates a morphism \(r_1 \to r_2\) in
  the implicit case category.
\item
  \textbf{Saturation}: Close the extracted morphism set under
  composition, identity, and the enriched weight constraints of
  \autoref{sec:enriched-categories}. Compute the empirical hom-values as
  normalized co-occurrence frequencies.
\item
  \textbf{Classification}: Construct the classifying topos
  \(\mathcal{E}_\mathbb{T}\) from the saturated theory---the canonical
  topos whose models are exactly the case-assignment patterns attested
  in the corpus. The topos-theoretic invariants (sort count, arity
  distribution, axiom count) combined with the magnitude and magnitude
  homology invariants of the enriched theory
  (\autoref{sec:enriched-categories}) provide a multi-dimensional
  fingerprint of the language's case system.
\item
  \textbf{Identification}: Compare the fingerprint against the Morita
  equivalence classes of known alignment types. If the classifying topos
  matches an existing class, the language's alignment type is
  identified; if not, the procedure has discovered a novel alignment
  pattern.
\end{enumerate}

This topos-theoretic learning procedure would be provably correct
(recovering the true theory in the limit) and maximally general (not
presupposing any particular alignment type). The magnitude invariant
from \autoref{sec:enriched-categories} enters at step 3 as a scalar
summary of the learned category's ``effective size,'' while magnitude
homology provides a finer-grained topological signature.
\end{frame}

\begin{frame}{Diagrammatic Implications of Inter-Theoretic Transfer}
\protect\phantomsection\label{diagrammatic-implications-of-inter-theoretic-transfer}
The bridge technique has a natural diagrammatic interpretation. Morita
equivalence between theories is witnessed by \emph{functorial
translations}---diagrams in the 2-category of toposes that commute up to
natural isomorphism. These diagrams serve the same cognitive function as
the commutative diagrams of \autoref{sec:introduction}: they make the
transfer of structure visible, allowing a researcher to verify at a
glance that a result proved in one framework genuinely applies in
another.

Manders {[}-@manders2008euclidean{]} observed that even in classical
mathematics, diagrams serve not merely as illustrations but as
\emph{inferential instruments} whose spatial properties encode
proof-relevant information. The topos-theoretic bridge diagrams extend
this observation to the meta-theoretical level: the commutative diagram
expressing Morita equivalence is itself a ``free ride'' inference,
automatically transferring any topos-invariant property from one theory
to another without requiring a case-by-case verification.

\textbf{Concrete bridge transfer.} Consider the \emph{commutativity of
transitive composition} proved in \(\mathbb{T}_{\text{log}}\) (the
type-logical framework): for types \(n, s\), the transitive derivation
\(n \cdot (n^r \cdot s \cdot n^l) \cdot n \to s\) yields the same
sentence type regardless of whether contractions proceed left-to-right
or right-to-left. Via the Morita equivalence chain, this result
transfers to \(\mathbb{T}_{\text{dist}}\) as a \emph{composition law}:
the sentence vector \(\overrightarrow{\text{Alice chases Bob}}\)
computed by contracting the subject tensor first (then the object)
equals the vector computed by contracting the object first (then the
subject)---a non-trivial commutativity property of the DisCoCat tensor
contraction that would require a separate linear-algebra proof without
the bridge. In \(\mathbb{T}_{\text{enr}}\), the same result becomes a
\emph{weight invariant}: the enriched hom-value of the composed morphism
NOM→ACC is independent of the intermediate case role through which
composition factors.
\end{frame}

\begin{frame}[fragile]{Computational Implementation of Topos-Theoretic
Bridges}
\protect\phantomsection\label{computational-implementation-of-topos-theoretic-bridges}
The topos-theoretic constructions developed above are not merely
abstract formalism---they are computationally implemented in our
\texttt{topos} module, which provides working Python implementations of
geometric theories, classifying toposes, Morita equivalence checking,
and bridge transfer.

\begin{block}{Geometric Theories from Case Categories}
\protect\phantomsection\label{geometric-theories-from-case-categories}
The \texttt{build\_typological\_theory()} function constructs a
geometric theory \(\mathbb{T}_{\text{typ}}\) from a
\texttt{CaseCategory} by extracting:

\begin{itemize}
\tightlist
\item
  \textbf{Sorts}: the case roles (objects of the category)
\item
  \textbf{Function symbols}: the morphisms with their source/target
  pairs
\item
  \textbf{Axioms}: identity morphism existence, composition closure, and
  alignment constraints
\end{itemize}

For the standard 8-case category, this yields a theory with 8 sorts and
approximately 15 function symbols. The mineral 3-case category produces
a theory with 3 sorts and 5 function symbols. Our
\texttt{build\_enriched\_theory()} function further annotates the
geometric theory with \([0,1]\)-valued hom weights from the enriched
structure of \autoref{sec:enriched-categories}.
\end{block}

\begin{block}{Classifying Toposes and Morita Equivalence Verification}
\protect\phantomsection\label{classifying-toposes-and-morita-equivalence-verification}
The \texttt{ClassifyingTopos} class computes topological
invariants---number of sorts, function arity distribution, and axiom
count---that characterize the ``logical shape'' of a theory. Two
theories are \textbf{Morita equivalent} when their classifying toposes
share the same invariant profile:

\begin{Shaded}
\begin{Highlighting}[]
\NormalTok{T\_std }\OperatorTok{=}\NormalTok{ build\_typological\_theory(standard\_case\_category())}
\NormalTok{T\_min }\OperatorTok{=}\NormalTok{ build\_typological\_theory(minimal\_case\_category())}
\NormalTok{equivalent }\OperatorTok{=}\NormalTok{ check\_morita\_equivalence(T\_std, T\_min)}
\CommentTok{\# False: 8{-}sort and 3{-}sort theories have different invariants}
\end{Highlighting}
\end{Shaded}

The \texttt{bridge\_transfer()} function implements the transfer
mechanism: given two Morita-equivalent theories, it constructs the
functorial translation that carries properties (alignment constraints,
composition laws) from one to the other. The transfer is blocked when
Morita equivalence fails, preventing unsound cross-theoretic
reasoning---a computational enforcement of the mathematical constraint.
\end{block}
\end{frame}

\end{document}
